% 第 32 章　信息、生命与熵
\chapter{信息、生命与熵}

在前面几章里，我们已经把热力学第二定律从“打碎的杯子不会自己复原”（第 29 章）一路推进到它的统计本质（第 31 章）：孤立系统的熵几乎总是增加，因为高熵的微观态数目以压倒性的优势多。这条定律听起来像是给全宇宙签了一张单程票。可是环顾四周，你会看到许多“逆行”的景象：冰箱把冷热分开，种子长成大树，细胞把散乱的小分子组装成精密机器。它们是在违法吗？本章要回答这个问题，而且答案会迫使我们承认一件更深刻的事：信息不是飘在物理世界之外的抽象符号，它寄居在物质里，记录它、复制它、擦除它，都是不折不扣的物理过程，都要遵守能量与熵的账本。

\step{反常识开场}

先看两个你天天见到、却很少放在一起想的场景。

场景一：夏天你把室温的饮料放进冰箱。几个小时后，冰箱内部冷热分明——冷冻室零下十几度，冷藏室四度，厨房里却微微变热。局部的“整齐”在增加：原来混在一起的分子运动，被人为地分成了“很安静的”和“很闹腾的”两群。

场景二：院子里一棵树苗，十年后长成大树。它把空气中弥散的二氧化碳、土壤里稀释的水和矿物质，装配成纤维排列整齐的树干、脉络分明的叶片。小分子的“乌合之众”变成了高度有序的结构。这棵树每时每刻都在制造秩序，而且理直气壮。

场景三其实就在你身上：此刻你的体温稳定在三十七度上下，你的细胞在不停地修复损伤、维持着膜两侧悬殊的离子浓度差、把氨基酸按精确的顺序串成蛋白质。只要活着，你就是一个不断抵抗“均匀化”的装置。

第 30 章说，孤立系统的熵不会减少。冰箱内部加一棵大树再加一个你，看起来都在“减熵”。难道第二定律有后门？如果有，我们是不是也能造一台只从海水吸热就永远运转的机器？

还有一个更隐蔽的问题，藏在你口袋里：删掉手机里 1 GB 的照片，手机没有变轻，照片里的“信息”去了哪里？彻底的“忘记”，在物理上是免费的吗？这一章结尾，你会发现这个问题的答案不仅解释了冰箱和大树，还降伏了一只纠缠了物理学家一百多年的小妖怪。

\step{先猜一猜}

在继续往下读之前，请把你的直觉写下来。对下面三个说法，猜“对”还是“错”：

\begin{enumerate}[leftmargin=2em]
\item 冰箱是在局部违反热力学第二定律，只是规模很小所以没关系。
\item 一株植物长大时变“有序”了，所以生命至少是第二定律的一个例外。
\item 理论上，把硬盘里 1 比特信息彻底擦除，至少要付出一份确定数量的能量，变成热散掉。
\end{enumerate}

再猜一个数量级问题：如果第 3 题是对的，擦除 1 比特的最小热代价，大约相当于下列哪一项？（a）烧开一壶水；（b）一粒尘埃从桌面落到地面的动能；（c）远远小于一粒尘埃落地的动能。先别翻答案，也别查资料。直觉越明确，待会儿的冲击越清楚；猜错不丢人，本章里猜错的都是大物理学家。

\step{动手实验}

\begin{experiment}[两张照片与两段文字]
这个实验不需要任何仪器，只要你的手机和电脑。

\textbf{第一步}：用手机在同一光线、同一分辨率下拍两张照片。第一张对着一面干净的白墙，第二张对着一张堆满杂物、颜色杂乱的书桌。查看两个文件的大小。

\textbf{第二步}：新建一个文本文件，输入一整屏“啊啊啊啊啊……”，另存为一份；再用键盘随机乱敲一整屏字符，另存为另一份。用压缩软件（ZIP 格式即可）分别压缩它们，比较压缩前后的大小。

\textbf{你会看到}：白墙照片可能只有几十 KB，杂乱书桌的照片却有几 MB，相差上百倍；“啊啊啊”文本被压缩到几乎为零，而随机乱敲的文本几乎压不动。

\textbf{这说明什么}：文件的大小——也就是记录它所需要的信息量——不取决于内容“有没有意义”，而取决于它有多\textbf{不可预测}。白墙完全可以预测（下一个像素还是白的），所以几乎不含信息；随机敲出的字符每一个都出乎意料，所以占满空间。注意，这和日常语言正好拧着：日常说“这张照片信息量很大”是指它内容丰富，而物理意义上信息量最大的，恰恰是看起来最“乱”的那一张。这是本章第一个让直觉摔跤的地方，请记住这种别扭感。

\textbf{附加观察}：摸摸运行中的冰箱后背或侧面，它是热的。再留意一下：冰箱运行时的嗡嗡声，意味着电能正在变成机械功。记住这个温度和这个声音，它们是本章的关键物证。
\end{experiment}

\step{旧模型遇到麻烦}

\textbf{第一个麻烦：“熵就是混乱程度”这句口头禅不够用。}我们在第 29 章就提醒过，不能把熵简单定义为“混乱程度”，现在你能看清为什么了。水在冰箱里结成冰，分子排成规整的晶格，按“混乱”的说法熵应该减小，可结冰明明是在零下的环境里自然发生的事。再看洗牌：一副按花色排好的新牌洗上七次，谁都同意“变乱了”，但牌面顺序的数目其实一直没变——$52!$ 种排列，洗牌前是其中一种，洗牌后还是其中一种。“乱”与“不乱”的差别不在排列的数目，而在我们对排法的分类方式。更根本的问题是：“混乱”无法测量——同一杯咖啡加牛奶，美术生看到的是好看的纹理，化学家看到的是均匀的溶液，物理学家却需要一个谁测都一样的数。熵的正确定义不靠观感，靠数数：一个宏观状态对应多少种微观排布，数目取对数，就是熵（第 31 章）。可数数这件事本身，立刻引出一个新角色——谁来数？

\textbf{第二个麻烦，也是最深的麻烦：麦克斯韦妖。}1867 年，麦克斯韦在给友人的信里提出一个思想实验。想象一个盒子被隔板分成左右两室，装满温度均匀的气体。隔板上有一扇无摩擦的小门，门口守着一个小东西——后人叫它“妖”。它的工作很简单：看见飞得快的分子从右边过来就开门放行，慢的就关门；反过来，只让慢分子从左边溜回右边。久而久之，左边越来越热，右边越来越冷，温差凭空出现。这个温差可以推动热机对外做功，而妖自己似乎什么能量都没花——它只是“看了看”，“记了记”，再开开关关一扇不费劲的门。

\begin{figure}[htbp]
\centering
\begin{tikzpicture}[scale=0.9]
  \draw[thick] (0,0) rectangle (8,3);
  \draw[thick] (4,0) -- (4,1.1);
  \draw[thick] (4,1.9) -- (4,3);
  \draw[thick,fill=gray!20] (4,1.1) rectangle (4.25,1.9);
  \node at (4.6,2.15) {\small 小门};
  % 妖
  \draw (3.55,2.5) circle (0.18);
  \draw (3.55,2.32) -- (3.55,2.0);
  \draw (3.55,2.2) -- (3.3,2.05);
  \draw (3.55,2.2) -- (3.85,2.05);
  \node at (3.0,2.6) {\small 妖};
  % 快分子
  \draw[-{Stealth[length=2.5mm]},thick] (6.3,1.5) -- (4.45,1.5);
  \fill (6.3,1.5) circle (0.09);
  \node at (6.3,1.9) {\small 快分子};
  % 慢分子
  \draw[-{Stealth[length=2.5mm]},thick] (1.7,0.8) -- (3.75,0.8);
  \fill (1.7,0.8) circle (0.09);
  \node at (1.7,0.45) {\small 慢分子};
  \node at (2,3.25) {\small 左边：越来越热};
  \node at (6,3.25) {\small 右边：越来越冷};
\end{tikzpicture}
\caption{麦克斯韦妖：一个“只看不做功”的守门员，似乎能不花代价地制造温差。}
\end{figure}

如果这套操作合法，第二定律就破产了：盒子和妖组成的系统孤立着，熵却减少了。一百多年里，物理学家提出了好几代“降妖方案”，每一代都被证明打偏了靶子。最早的方案针对门：斯莫卢霍夫斯基在 1912 年前后指出，一扇轻到能被分子撞开的小门，自己也会在热运动下乱晃乱开，快分子慢分子一起漏，妖的筛选自动失效。但这只打败了“自动门”，没有打败“妖”——一个会判断、会主动开门的装置不受此限。20 世纪 50 年代的方案针对眼睛：布里渊论证，黑暗的盒子里妖要看清分子就得开灯，光子打在分子上扰乱系统，耗费不小。可这个论证也不彻底——理论上可以设计不靠光照的测量方式。要害始终没有被击中：妖必须\textbf{获取并记住}每个分子是快是慢。观察和记忆，到底是不是免费的？这个问题从 1867 年一直悬到 20 世纪后半叶。现代分子机器的存在让问题更加紧迫：细胞里真实运转的分子马达比妖还小，“在分子尺度上精细操作”本身没有原则障碍，那么第二定律的防线究竟筑在哪里？

\step{发明新概念}

回答这个问题的过程，逼出了 20 世纪最重要的新概念之一：\textbf{信息是物理的}。

先弄清“信息”本身。1948 年，香农建立了信息论。他的出发点冷酷而有效：忘掉“意义”，只数“可能性”。一个问题的答案如果只有两种等概率的选项（硬币正还是反），告诉你答案，就给了你 1 \textbf{比特}（bit）的信息；四种等概率选项，答案是 2 比特；一枚均匀骰子的点数，约 2.58 比特。可能性越多、越出人意料的结果，包含的信息越多。你在动手实验里已经亲手验证了这个定义：白墙照片“下一个像素还是白的”，等于没有新可能性，所以信息量极小。这套理论不管消息“重不重要”——“明天天气晴”和“明天下金子”只要概率相同，信息量就相同；它只负责称量“不确定性的减少”，像一台不看货色只称重的秤。

这套秤在今天的基础设施里无处不在。手机拍照生成的 JPEG 文件、聊天软件里的 ZIP 压缩包，用的都是“可预测性就是可压缩性”这一条；二维码被咖啡渍弄脏一块照样能扫出来，旅行者号探测器从两百多亿公里外发回的微弱信号照样能读出照片，靠的是纠错码——故意加入有规律的冗余，让接收方能从残缺的信息里还原原样。这些都是信息论的工程化身。但香农的理论本身不涉及任何原子：它完全可以在纸上用笔算完成。真正把它拽进物理世界的，是接下来这个追问：信息存放在哪里？

存放在载体上。一条信息要存在，就必须写在某个物理系统里：纸上的墨水颗粒、硬盘的磁畴取向、闪存浮栅里的电荷、DNA 的碱基序列、脑细胞之间的连接强度。没有载体就没有信息——“信息是物理的”，这是兰道尔的名言。于是记录、复制、擦除信息，统统是对原子做的物理操作，统统要进热力学账本。

本案的关键证词，分三段。

\textbf{第一段，1929 年，西拉德}把麦克斯韦妖简化到不能再简：盒子里只放一个分子。妖测量分子在左还是在右（获得 1 比特信息），然后在对应一侧插入活塞，让分子等温膨胀对外做功，取出能量，再拔回活塞，循环往复。整台机器似乎只靠“知道”就从单一热源里榨出了功。西拉德指出：漏洞不在别处，正在“获取信息”这一步——测量不可能白白发生。他第一次把“1 比特信息”和一份确定的热力学量挂上了钩。

\textbf{第二段，1961 年，兰道尔}给出了精确得多的指控。他研究的不是“看”，而是“忘”。擦除 1 比特——把“可能是 0 也可能是 1”的载体强行重置成确定的 0——是把两个可能状态压缩成一个。这一步在逻辑上不可逆：面对重置后的那个 0，你永远无法反推它原来是 0 还是 1。兰道尔证明，逻辑不可逆必然伴随着物理上的放热：每擦除 1 比特，至少要向环境放出一份确定的热量。这就是兰道尔原理。反过来，逻辑上可逆的操作——比如把信息从一处复制到另一处，原样保留——原则上可以做到不放热。记录可以轻巧，遗忘却要付费；这条不对称是整章的枢纽。

\textbf{第三段，1982 年，贝内特}把案子结了。他证明：只要操作足够慢、足够小心，测量本身可以做成可逆的，不必放热——妖“看”分子这一下可以是免费的。但妖的记忆容量有限，它不可能无限地记下去；要循环工作，就必须清空记忆给下一次测量腾地方，而清空就是擦除，擦除就要放热。放出的热恰好把妖从温差里赚来的好处全部抵消，一分不多，一分不少。妖不是被“看”累死的，是被“忘”累死的。一百多年的悬案，靠“信息擦除要付热账”这一条新定律告破。第二定律的防线，原来筑在记忆体的垃圾桶旁。

这个视角顺手解开了另一个日常之谜。把热咖啡放在桌上，奶花的形状、你拉花时手腕的抖动，这些“信息”半小时后就永远消失了——不是因为它们被擦掉，而是因为它们被搅拌进了房间的空气里：奶花瓦解时，分子的排布方式被冲进了数量大得无法追踪的环境自由度中。从一杯凉咖啡反推当初的奶花，等于要从几十亿亿个空气分子的轨迹里还原一段历史，这在实践中和信息被彻底擦除没有区别。熵增的每一步，都是信息从“我们够得着的地方”流向“我们够不着的地方”的过程。过去之所以确定、未来之所以模糊，根源就在这趟单向的搬运里。

这里值得停下来声明一下本章反复使用的比喻。我们说熵和信息要“记账”，它像会计账本：每一笔都要平衡，局部科目的盈余必然对应别处更大的支出。但它不像账本的地方在于：账本上记的数字不是钱本身——熵不是某种流体物质，信息也不是一种新能量；它们是状态的计数方式，是\textbf{描述}系统的语言，只是这种语言恰好被物理定律严格约束着。把“熵”当成一瓶可以倒来倒去的灰色液体，是这个比喻最常见的翻车方式。

还要分清两个容易混淆的词。\textbf{信息熵}（香农熵）数的是“消息的可能性”，单位是比特；\textbf{热力学熵}数的是“分子的微观排布”，单位是焦耳每开尔文。它们刻画的对象不同，平时各用各的，不能随意互换。但它们在同一张桌子上结账：当信息被写进物理载体、又被物理地擦除时，两种熵就被换算关系拴在了一起——这正是兰道尔原理的内容，也是下一节要用数学写清楚的东西。

\step{一句话模型}

\begin{oneline}
信息寄居在物质里，擦除它必须放热；生命和冰箱都不违反第二定律——它们只是熵洪流中一个临时的小旋涡，把更多的熵推给了环境。
\end{oneline}

\step{数学翻译器}

现在把上面的故事翻译成式子。

\textbf{第一式：信息量。}一个消息有 $n$ 种可能的结果，第 $i$ 种出现的概率是 $p_i$，那么它的香农熵定义为
\[
H=-\sum_{i=1}^{n}p_i\log_2 p_i \qquad \text{（单位：比特）}.
\]
这个式子在说什么：每一项是“该结果的概率”乘以“该结果的意外程度”（概率越小，$\log_2$ 越大，越意外），加起来就是平均意外程度。为什么这样写：它是唯一同时满足“确定事件信息量为零”“可能性越多信息量越大”“两次独立选择的信息量可以相加”这三条常识要求的函数。立刻用特殊情况检查它：

\begin{itemize}[leftmargin=2em]
\item 结果完全确定（$p_1=1$）：$H=-1\times\log_2 1=0$。意料之中，没有信息量。白墙照片就是这种极限。
\item $n$ 个结果等概率：$H=-n\times\frac1n\log_2\frac1n=\log_2 n$。硬币给出 $\log_2 2=1$ 比特，骰子给出 $\log_2 6\approx2.58$ 比特，符合预期。
\item 某个 $p_i\to 0$：$p_i\log_2 p_i\to 0$，几乎不可能的事件对平均信息量几乎没有贡献，公式依然光滑。
\end{itemize}

\textbf{第二式：遗忘的价目表。}在温度 $T$ 的环境里擦除 1 比特信息，向环境放出的热量至少为
\[
Q_{\min}=k_{\mathrm{B}}T\ln 2 ,
\]
其中 $k_{\mathrm{B}}=1.38\times10^{-23}\ \mathrm{J/K}$ 是玻尔兹曼常数。为什么恰好是 $k_{\mathrm{B}}T\ln 2$：擦除 1 比特是把载体的两个可能状态压成一个，载体丢掉的那份自由度必须由环境接收，环境熵至少增加 $k_{\mathrm{B}}\ln 2$，乘以温度就是最少放出的热。检查特殊情况：$T\to 0$ 时代价趋于零——看似免费，但第 30 章的第三定律说绝对零度只能逼近不能到达，所以“零成本的遗忘”是极限而不是现实；若信息本来就是确定的（不存在“可能是 0 也可能是 1”的含糊），重置它不算擦除，不收费；方向反转也成立——把一份比特\textbf{复制}到空白载体上，逻辑可逆，原则上可以不付热账。收费的从来不是信息本身，而是“让多种可能塌缩成一种”的不可逆动作。

代入真实数字感受一下。室温 $T\approx300\ \mathrm{K}$：
\[
Q_{\min}\approx1.38\times10^{-23}\times300\times0.693\approx2.9\times10^{-21}\ \mathrm{J},
\]
约合 $0.018\ \mathrm{eV}$。把整部手机 1 GB（约 $8\times10^9$ 比特）的照片全部彻底擦除，理论最小发热量只有约 $2\times10^{-11}$ 焦耳。一粒尘埃落地释放的能量都比它大上万倍。所以“删照片”在你手机里产生的那点热，几乎全是芯片效率低下造成的，离理论极限还很远——今天的芯片每个逻辑操作的耗散仍比兰道尔极限高出成千上万倍。但极限本身是真的：2012 年，实验物理学家用激光镊子控制一颗微米级胶体颗粒在双势阱中充当 1 比特，缓慢地执行擦除，实测放出的热量逼近了 $k_{\mathrm{B}}T\ln 2$，与理论吻合。妖的账单，被实验亲自签收。

这张价目表也是一道工程天花板。半个世纪以来，芯片上晶体管越造越小、每个操作的能耗一路下降，但照这个趋势，总有一天会撞上兰道尔极限这堵墙——到那时，继续降低功耗只有两条路：要么降低运行温度（代价高昂），要么改变计算的逻辑结构，让每一步都可逆、不擦除。后面这条路叫可逆计算，今天还只是实验室里的探索，但它是信息论给芯片工业提前写下的一张通知单。

眼前的能耗账同样可观。全世界的数据中心每年消耗的电力已达数千太瓦时，占全球用电量的百分之一到百分之二，相当于一个中等工业国家的全部用电；其中相当大一部分花在散热上——风扇和空调昼夜不停地把芯片擦除信息、翻转比特产生的热搬运到大气里。你每一次搜索、每一段视频，背后都有一小股实实在在的熵被排进环境。信息社会并不“轻”，它的重量写在电厂的煤耗和冷却塔的蒸汽里。

\textbf{第三笔账：你一天排出多少熵。}一个成年人一天从食物获得约 2000 千卡，合 $8.4\times10^6$ 焦耳。这些能量最终几乎都以热的形式排给环境。设环境温度约 $295\ \mathrm{K}$，粗略估计你每天向外输出的熵是
\[
\Delta S_{\text{排}}\approx\frac{8.4\times10^6\ \mathrm{J}}{295\ \mathrm{K}}\approx2.8\times10^4\ \mathrm{J/K}.
\]
你体内维持的那点秩序——整齐折叠的蛋白质、浓度精确的细胞液、排列成编码的 DNA——所对应的熵“亏空”，和这三万焦耳每开尔文的支出相比只是零头。生命不是不向熵交税，而是交得又勤又多。

\begin{mathbridge}
\textbf{两种熵的换算与西拉德引擎。}第 31 章给出热力学熵的微观定义 $S=k_{\mathrm{B}}\ln\Omega$，其中 $\Omega$ 是微观态数目。若 $\Omega$ 个微观态等概率，香农熵 $H=\log_2\Omega$，于是
\[
S=k_{\mathrm{B}}\ln\Omega=(k_{\mathrm{B}}\ln2)\,H.
\]
也就是说：1 比特信息熵，写在物理载体上，恰好对应 $k_{\mathrm{B}}\ln2$ 的热力学熵。更一般地，把香农熵的求和式乘上 $k_{\mathrm{B}}\ln2$，就得到统计物理中的吉布斯熵公式——两套语言在这里完成了对译。现在回头看西拉德的单分子引擎：分子等温膨胀，体积加倍，微观态数翻倍，从热源吸热并对外做功
\[
W=k_{\mathrm{B}}T\ln 2 .
\]
妖每循环一次赚了 $k_{\mathrm{B}}T\ln2$ 的功，但它的记忆里多了一条必须清掉的记录；清空这条记录至少放出 $Q_{\min}=k_{\mathrm{B}}T\ln2$ 的热。$W-Q_{\min}=0$：赚多少，赔多少，第二定律分文未损。$\ln2$ 这个小小的数字，站在等式两边，把信息论和热力学缝在了一起。
\end{mathbridge}

\begin{challenge}
\textbf{小系统的例外、可逆计算与黑洞的熵。}三件事供想深挖的读者。

其一，兰道尔原理和第二定律一样是统计性的。近几十年被统称为“涨落定理”的一批严格结果表明：在足够小的系统、足够短的时间里，熵瞬时减少是允许的，只是其概率随偏离程度指数衰减。微米尺度、皮秒尺度的“违章”已被实验观测到；尺度放大到宏观，违章概率小到宇宙年龄里都不会发生一次。

其二，既然只有擦除必须放热，那么一台\textbf{从不擦除}的计算机——每一步都保留足够信息以便原路倒回——原则上可以把耗散压到任意低。这种“可逆计算”已有严格的逻辑门设计（如弗雷德金门、托福利门），是低功耗芯片和量子计算的一条重要思想源头。

其三，熵与信息的联姻在最极端的地方开出了最奇怪的花：黑洞。20 世纪 70 年代，贝肯斯坦和霍金发现黑洞的熵正比于其视界面积，
\[
S_{\mathrm{BH}}=\frac{k_{\mathrm{B}}c^{3}}{4G\hbar}\,A,
\]
一个太阳质量黑洞的熵高达 $10^{54}\ \mathrm{J/K}$ 量级。面积而不是体积出现在公式里，暗示“一个空间区域能装下的信息存在上限”，这条线索（全息原理）如今是量子引力研究的核心之一，我们将在第 41 章重新遇见它。
\end{challenge}

\step{模型的边界}

\textbf{边界一：第二定律是统计定律，不是铁律，但尺度是生死线。}对几个粒子组成的系统，熵自发减少的涨落俯拾皆是——布朗运动中一颗花粉颗粒被水分子撞得突然加速，就是一次微观的“熵减少”，毫不神秘，也不违反任何定律；你甚至可以做个粗略的类比：三个人猜拳，三人恰好出成一模一样的局面时有出现；三亿人猜拳，要全体出成一样，宇宙灭亡前都等不到一回。粒子数就是那三亿人的亿万倍。问题在于外推：$10^{23}$ 个粒子的宏观系统里，第二定律被违反的概率不是“小”，而是小到你写下这个概率所需的零比可见宇宙的原子还多。“几乎不可能”和“绝对禁止”的区别（第 30 章）在这里必须咬死：第二定律允许例外，但例外被囚禁在微观尺度。

\textbf{边界二：熵增不等于“宇宙里所有东西都会立刻变乱”。}这是本章必须拆掉的最后一个误解。第二定律管的是\textbf{孤立系统的总账}，不禁止任何局部的、暂时的有序化，只要它向环境输出的熵更多。水结成冰、熔岩凝成晶体、胚胎长成婴儿，都是合法的局部减熵。更微妙的例子来自引力：一团均匀弥散的气体云，对引力系统来说其实是\textbf{低熵}状态；它在引力下收缩、碎裂、点燃成恒星、聚集成星系，看起来结构越来越丰富，熵却在一路增加。引力的存在让“均匀”与“高熵”脱了钩——这正是宇宙既没有立刻热寂、又能长出星系和生命的原因。所谓“热寂”是遥远的渐近图景，不是此刻正在上演的剧情。

\textbf{边界三：别把本章的概念借去乱用。}两个最常见的误用。一是把热力学熵当成“混乱”的同义词套到社会、人生、办公桌上去——那是修辞，不是物理，社会系统不数微观态。二是把信息熵和热力学熵当成一回事——前者数消息的可能性，后者数分子的排布，只有在“信息被物理地记录或擦除”这样的具体过程里，两者才通过 $k_{\mathrm{B}}\ln2$ 换算。看到“某某系统的熵在增加”这类流行说法，先问一句：他数的是什么态？

\textbf{边界四：生命不“对抗”熵。}流行的诗意说法是“生命以负熵为食”，这句话有物理内核（薛定谔 1944 年的表述），但容易被误读成生命在抵抗第二定律。准确的说法平淡得多：生命是一个敞开的通道，低熵的能量和物质流进来，高熵的废物和热量流出去，通道内部维持精巧的结构，全靠流量而不是豁免。流量一停（死亡），通道立刻崩塌回平衡态。同一原理也写在城市的地图上：盛夏里每一台空调都在给室内制造凉爽的秩序，同时把更多的热连同电费转化的能量吐到街上——整座城市因此比郊区更热。局部的舒适和大气的升温，是同一张账单的正面和背面。

\step{回头解释开场现象}

现在结算开场的三笔悬念。

\textbf{冰箱}：冰箱不是把“热消灭掉”，而是用压缩机消耗电功，把热量从冷端搬运到热端。你从后背摸到的那股热，等于从箱内抽走的热\textbf{加上}电费转化来的功。排给厨房的热量总是大于从箱内吸走的热量，高温端放出的熵超过低温端吸收的熵，把冰箱和厨房算作一个整体，熵只增不减。局部的“冷热分明”，是花电费买来的，账是平的。

\textbf{大树}：植物收到的是太阳光——每个可见光光子能量高（约 2 eV）、数目少，携带的熵低；它向环境放出的是红外热辐射——同样的能量分摊给几十倍数量的低能光子（约 0.05 eV），熵高得多。一进一出，能量平衡，熵却净增。整套光合作用只截留这股东风的一小部分来合成有序的组织，其余全部以更高熵的形式退还。把镜头拉远到整个地球，账本最漂亮：地球平均每平方米从太阳吸收约 240 瓦，对应光子温度约 5800 K；又以约 255 K 的温度向太空辐射同样多的能量。流出的熵每秒钟每平方米约比流入的多出 0.9 焦耳每开尔文——地球上全部的风、河流、森林和你，都是这股巨大熵流冲刷出来的次级结构。连桌上那杯慢慢变凉的热咖啡也是其中一员：它把热量连同奶花的信息一起交给房间，参与的正是同一场单向搬运。

\begin{figure}[htbp]
\centering
\begin{tikzpicture}[scale=1.0]
  % 太阳
  \draw[fill=yellow!40] (-4.5,0) circle (0.6);
  \node at (-4.5,-1.1) {\small 太阳};
  % 地球
  \draw[fill=blue!15] (2.5,0) circle (0.8);
  \node at (2.5,0) {\small 地球};
  % 入射
  \draw[-{Stealth[length=3mm]},very thick] (-3.7,0) -- (1.5,0);
  \node[above] at (-1.1,0.15) {\small 少量高能光子：$240\ \mathrm{W/m^2}$，约 $5800\ \mathrm{K}$};
  \node[below] at (-1.1,-0.15) {\small 低熵流};
  % 出射（多箭头）
  \foreach \y in {0.5,0.17,-0.17,-0.5}{
    \draw[-{Stealth[length=2.5mm]},thick,decorate,decoration={snake,amplitude=0.6mm}] (3.3,\y) -- (5.6,\y);}
  \node[right] at (5.6,0.3) {\small 大量低能红外光子，约 $255\ \mathrm{K}$};
  \node[right] at (5.6,-0.35) {\small 高熵流，能量相同、熵更多};
\end{tikzpicture}
\caption{地球的熵账本：能量收支平衡，熵却大额净出口。生命是这笔差额的受益者。}
\end{figure}

\textbf{删照片}：照片的信息记录在闪存单元的电荷里。删掉它们，若真的执行物理上的彻底重置，每比特至少要付 $k_{\mathrm{B}}T\ln2$ 的热。你的猜测现在可以对答案了：第 1 题错，第 2 题错，第 3 题对，数量级是（c）——远远小于一粒尘埃落地的动能。世界上没有免费的遗忘，只有便宜到感觉不到的遗忘。

\step{脑洞实验与挑战题}

本章的故事还会在后面的章节里长出新枝：当“比特”换成遵守量子规则的“量子比特”，信息的记录、复制和擦除会呈现全新的限制（比如量子态无法被完美复制），那是第 47、48 章的题材；而“测量究竟怎样改变一个系统”这个问题，将在第 45 章被重新审视。

\begin{thought}[给妖配一块无限大的硬盘]
贝内特的结论是：妖败在记忆有限，必须擦除。那如果给它一块\textbf{永不装满}的记忆体呢？它可以永远只记不删，第二定律是不是就绕过去了？

想一想这台“永动机”的终局：盒子里的温差确实被它榨成了功，但它的记忆体在不停膨胀——每一个分子的来龙去脉都写在里面。热力学第二定律并没有被违反，而是被迫改写成了更深刻的形式：熵不能消灭，只能\textbf{转移和寄存}。妖等于把气体的熵搬进了自己的记忆，并把房租（最终擦除的热）无限期拖欠下去。可宇宙里没有真正的无限硬盘——任何物理载体的容量都有限（黑洞熵公式甚至给出了单位面积的上限），账单到期的那一天，连本带利分文不少。熵是个极有耐心的债主。
\end{thought}

\begin{puzzles}
\item \textbf{为什么没有万能压缩算法？}有人宣称发明了一种算法，能把\textbf{任何}文件都至少压短 1 比特。不用看代码，证明他一定是骗子。（提示：长度为 $n$ 比特的文件共有 $2^n$ 个，而比它短的编码一共只有 $1+2+\dots+2^{n-1}=2^n-1$ 个——编码不够用，必有文件无处可压。这就是动手实验中“随机文本压不动”的根本原因：大多数字符串已经接近不可压缩。）

\item \textbf{打开冰箱门能当空调用吗？}盛夏的厨房里，有人敞开冰箱门想给房间降温。房间温度最终会升、降还是不变？如果把冰箱和房间严格绝热起来通电运行一天，结果又如何？（提示：冰箱从箱内吸走多少热，就往后背排出多少热，还要加上全部电费转化成的热。敞开门的冰箱是一台纯粹的电暖器，而且还很吵。）

\item \textbf{偏一点，反而省空间。}公平硬币连抛 1000 次的记录需要多少比特来存储？如果硬币被做了手脚，正面概率变成 $51\%$，同样的 1000 次记录需要的比特数会变多还是变少？（提示：用 $H=-p\log_2 p-(1-p)\log_2(1-p)$ 计算单次信息量。$51\%$ 时 $H\approx0.9997$ 比特，略小于 1。可预测性就是可压缩性，哪怕只偏一点点。）

\item \textbf{妖的最后一招：只看不记。}假设妖有一种“过目即忘”的超能力：测量分子的瞬间不留下任何记录，仅凭直觉开门。它能逃过兰道尔的账单吗？（提示：不能。“不留下记录的测量”等于没有测量——开门的决策本身就是一份物理记录，哪怕只存在一瞬，重置它的动作同样要结账。想占第二定律的便宜，连“忘了自己看过”这条路都被堵死了。）
\end{puzzles}
